%% Conclusion generale — 1 page
\chapter*{Conclusion générale}
\addcontentsline{toc}{chapter}{Conclusion générale}

Le projet \emph{BMS Engineering Assistant} a atteint l'ensemble de ses
objectifs. Deux livrables complémentaires ont été produits et intégrés~:
un BMS \textsc{Simulink} complet (mod\`{e}le ECM~2RC, FSM de protection,
équilibrage, pronostic LSTM) validé sur soixante-dix trajets mesurés~;
et un pipeline logiciel automatisé qui pilote la campagne de validation,
en analyse les résultats par un moteur de règles YAML, puis synthétise un
rapport HTML exploité par un LLM.

Les résultats clés obtenus sont les suivants~: toutes les 86~exigences
couvrent à 100~\% après intégration~; le temps d'exécution WCET PIL
mesuré (1000{,}6~µs) représente moins de 11~\% du budget de 10~ms~;
le prédicteur LSTM détecte les trois types de défauts avec une latence
inférieure à 250~ms et zéro faux-positif sur le dataset de référence~;
la campagne complète s'exécute en moins de 15~minutes à froid, et en
moins de 2~minutes avec les caches actifs.

Trois axes principaux d'évolution se dessinent. L'extension du catalogue
de défauts aux scénarios combinés (déséquilibre simultané à un emballement
thermique) constitue la priorité de court terme. Le portage à une cible
automobile certifiable (\emph{Infineon AURIX TC4xx}~/ \emph{NXP S32K3})
permettra d'apprécier la robustesse de la chaîne \emph{Embedded Coder}
face aux contraintes ASIL-B. Enfin, la substitution du service Groq par
un modèle local (\textsc{vLLM}~/ \textsc{Ollama}) est immédiatement
faisable et garantira l'indépendance vis-à-vis de tout fournisseur
exterieur pour une mise en production à l'échelle industrielle.
